\section{Setting and main results}\label{sec:setting}

\subsection{Graphs and standing assumptions}\label{sec:graphs}

Let $G=(V,E)$ be a graph with vertex set $V$ and edge set $E$. We write $x\sim y$ if $\{x,y\}\in E$,
$\dist$ for the graph distance, $B(x,r)$ for the closed ball of radius~$r$ around~$x$, and
$\Aut(G)$ for the group of graph automorphisms. For $A\subseteq V$ let $\partial_E A$ be the set of
edges with exactly one endpoint in~$A$, and let $\Nb[A]$ be the closed neighbourhood of~$A$, consisting
of $A$ and all vertices adjacent to a vertex of~$A$. The graph is \emph{transitive} if $\Aut(G)$ acts
transitively on~$V$, and \emph{nonamenable} if
\begin{equation*}
  \iota \coloneqq \inf\Bigl\{\frac{\abs{\partial_E A}}{\abs{A}} \colon A\subseteq V \text{ finite},\ A\neq\emptyset\Bigr\} > 0\,.
\end{equation*}
For graphs of bounded degree this is equivalent to positivity of the vertex-isoperimetric constant,
because every vertex of the outer vertex boundary of~$A$ is incident to at least one and at most
boundedly many edges of~$\partial_E A$.

Throughout these notes $G$ satisfies the following standing assumption.

\begin{asslist}
  \item\label{ass:G} $G=(V,E)$ is an infinite, connected, locally finite, simple graph, $\Aut(G)$ acts
  transitively on~$V$, and $G$ is nonamenable.
\end{asslist}

\noindent A transitive graph is regular; we write $d$ for its degree and fix a root vertex $o\in V$.
From \cref{sec:strategy} on, the main argument also uses the second assumption
\begin{asslist}[resume]
  \item\label{ass:U} the group $\Aut(G)$ is unimodular,
\end{asslist}
\noindent that is, $\abs{\operatorname{Stab}(x)y}=\abs{\operatorname{Stab}(y)x}$ for all $x,y\in V$, where
$\operatorname{Stab}(x)y$ is the orbit of $y$ under the stabiliser of~$x$ in $\Aut(G)$. Equivalently, the
mass-transport principle holds (\cref{sec:mtp}); every Cayley graph satisfies \cref{ass:U}. The
nonunimodular case is covered by a theorem of Hutchcroft (\cref{prop:nonunimodular}).

\subsection{Percolation, connection kernel and thresholds}\label{sec:percolation}

Let $\PP_p$ be the law of Bernoulli bond percolation with parameter $p\in[0,1]$ on~$G$, that is, the
product measure on $\{0,1\}^E$ under which each edge is \emph{open} ($\omega(e)=1$) with
probability~$p$, and let $\EE_p$ be its expectation. We write $x\conn y$ if $x$ and $y$ are joined
by an open path; paths of length zero are allowed, so $x\conn x$. The \emph{cluster} $C_x$ of~$x$ is
the set of vertices $y$ with $x\conn y$, and $\abs{C_x}$ is its number of vertices. For a vertex set
$\Lambda\subseteq V$, a \emph{domain}, we write $x\connin{\Lambda}y$ if $x$ and $y$ are joined by an open
path using only vertices of~$\Lambda$, and $C_x^\Lambda$ for the cluster of~$x$ in the subgraph induced
by~$\Lambda$. By convention $C_x^\Lambda=\emptyset$ if $x\notin\Lambda$, and every quantity that
requires a terminal outside its domain is zero. Let
\begin{equation*}
  \theta(p)\coloneqq\PP_p(\abs{C_o}=\infty)\,,\qquad
  p_c\coloneqq\inf\{p\colon\theta(p)>0\}\,,\qquad
  p_u\coloneqq\inf\{p\colon\PP_p(N_\infty=1)=1\}\,,
\end{equation*}
where $p$ ranges over $[0,1]$ and $N_\infty$ is the number of infinite clusters. In the \emph{standard monotone coupling} the
edges carry independent labels $U_e$ uniform on $[0,1]$, and an edge is open at parameter~$p$ if
$U_e\le p$; this realises all~$\PP_p$ on one probability space.

\paragraph{The connection kernel.}
A \emph{kernel} is a function $K\colon V\times V\to\R$. It is \emph{invariant} if $K(gx,gy)=K(x,y)$
for all $g\in\Aut(G)$ and $x,y\in V$. We let a kernel act on finitely supported functions by
$(Kf)(x)\coloneqq\sum_yK(x,y)f(y)$ and write $\norm{K}\in[0,\infty]$ for its operator norm on
$\ell^2(V)$, the Hilbert space of square-summable functions on~$V$ with counting measure; $K$ is
\emph{bounded} if $\norm{K}<\infty$. Products of nonnegative kernels,
$(KL)(x,y)\coloneqq\sum_zK(x,z)L(z,y)\in[0,\infty]$, are always defined. We write $\Adj$ for the
adjacency kernel, $\Adj(x,y)\coloneqq\1\{x\sim y\}$, and $I$ for the identity. The \emph{connection
kernel} and the \emph{operator threshold} are
\begin{equation*}
  T_p(x,y)\coloneqq\PP_p(x\conn y)\,,\qquad
  \ptwo\coloneqq\sup\{p\in[0,1]\colon \norm{T_p}<\infty\}\,.
\end{equation*}
More generally $T_p^\Lambda(x,y)\coloneqq\PP_p(x\connin{\Lambda}y)$. The \emph{susceptibility} and the
\emph{triangle diagram} are
\begin{equation*}
  \chi(p)\coloneqq\EE_p\abs{C_o}=\sum_{y\in V}T_p(o,y)\,,\qquad
  \nabla_p\coloneqq\sum_{y,z\in V}T_p(o,y)\,T_p(y,z)\,T_p(z,o)\,,
\end{equation*}
both with values in $[1,\infty]$. By transitivity they do not depend on the choice of~$o$.

\paragraph{Constants.}
Unless indicated otherwise, $c,C\in(0,\infty)$ denote constants that depend only on~$G$; their values
may change from line to line. Dependence on further parameters is shown by subscripts, as in~$C_\eta$.

\subsection{Main results}\label{sec:main-results}

\begin{theorem}[Critical boundedness and nonuniqueness]\label{thm:main}
Assume~\cref{ass:G}. Then
\begin{equation}\label{eq:main}
  \sup_{0\le p<p_c}\norm{T_p}=\norm{T_{p_c}}<\infty
  \qquad\text{and}\qquad
  p_c<\ptwo\le p_u\,.
\end{equation}
Moreover there are $p_c<p_1<p_2<1$ such that, in the standard monotone coupling, almost surely for
every $p\in[p_1,p_2]$ the open subgraph at parameter~$p$ has infinitely many infinite clusters.
\end{theorem}

The strict inequality $p_c<p_u$ is the bond-percolation case of the conjecture of Benjamini and
Schramm for transitive graphs.

\begin{corollary}[Cayley graphs]\label{cor:cayley}
Let $\Gamma$ be a finitely generated nonamenable group and $S\subseteq\Gamma$ a finite symmetric
generating set not containing the identity. Then the Cayley graph $\Cay(\Gamma,S)$, with vertex set
$\Gamma$ and edges $\{x,xs\}$ for $x\in\Gamma$ and $s\in S$, satisfies the conclusions of
\cref{thm:main,thm:consequences}.
\end{corollary}

\begin{proof}
The Cayley graph is simple, because the identity is not in~$S$ and edges are unordered pairs. It is
connected because $S$ generates~$\Gamma$, and locally finite because $S$ is finite. Left
multiplications are automorphisms, so it is transitive. By F{\o}lner's criterion a finitely generated
group is nonamenable if and only if one, equivalently every, of its Cayley graphs with respect to a
finite generating set is nonamenable (see~\cite[Chapter~6]{LP16}). % TODO: Referenz prüfen (exact locator in LP16)
Hence \cref{ass:G} holds, and \cref{thm:main,thm:consequences} apply.
\end{proof}

The operator bound also determines the critical behaviour. For a vertex $v$, the \emph{extrinsic} and
\emph{intrinsic radius} of its cluster are
\begin{equation*}
  \radext(C_v)\coloneqq\sup_{x\in C_v}\dist(v,x)\,,\qquad
  \radint(C_v)\coloneqq\sup_{x\in C_v}\dist_\omega(v,x)\,,
\end{equation*}
with values in $[0,\infty]$, where $\dist_\omega$ is the graph distance in the open subgraph. We say
that connection probabilities \emph{decay exponentially} at~$p$ if there are $A<\infty$ and $a>0$ with
$T_p(x,y)\le A\e^{-a\dist(x,y)}$ for all $x,y\in V$, and we let $\pexp$ be the supremum of all such~$p$.
We write $f\asymp g$ in a stated limit if $cg\le f\le Cg$ in a neighbourhood of that limit.

\begin{theorem}[Mean-field behaviour and exponential connectivity]\label{thm:consequences}
Assume~\cref{ass:G}. Then
\begin{equation}\label{eq:triangle-condition}
  \nabla_{p_c}\le\norm{T_{p_c}}^3<\infty\,,
\end{equation}
and, as $p\uparrow p_c$, $p\downarrow p_c$ and $n\to\infty$ respectively,
\begin{align}
  \chi(p)&\asymp(p_c-p)^{-1}\,,&
  \theta(p)&\asymp p-p_c\,,\label{eq:chi-theta}\\
  \PP_{p_c}(\abs{C_o}\ge n)&\asymp n^{-1/2}\,,&
  \PP_{p_c}\bigl(\radint(C_o)\ge n\bigr)\asymp\PP_{p_c}\bigl(\radext(C_o)\ge n\bigr)&\asymp n^{-1}\,.
  \label{eq:critical-tails}
\end{align}
Moreover $p_c<\ptwo\le\pexp$, and for every $p\in[0,\ptwo)$ there are $A_p<\infty$ and $a_p>0$ with
\begin{equation}\label{eq:exp-decay}
  T_p(x,y)\le A_p\,\e^{-a_p\dist(x,y)}\qquad\text{for all }x,y\in V\,.
\end{equation}
\end{theorem}

The bound~\eqref{eq:triangle-condition} is the \emph{triangle condition} of Aizenman and
Newman~\cite{AN84}. Together with $p_c<\pexp$ it verifies Conjectures~1.2 and~2.1 of~\cite{Hut20L2} for
transitive graphs. The asymptotics~\eqref{eq:chi-theta} and~\eqref{eq:critical-tails} are the
mean-field values of the critical exponents. They follow from known conditional results, whose
hypothesis is either the triangle condition or $p_c<\ptwo$; the new input is that these hypotheses hold
for every graph satisfying~\cref{ass:G}. The deduction is given in \cref{sec:consequences}.

\begin{remark}[Fixed parameters]\label{rem:fixed-p}
For each fixed $p\in(p_c,p_u)$ there are almost surely infinitely many infinite clusters. Indeed, on a
transitive graph $N_\infty$ is almost surely constant and belongs to $\{0,1,\infty\}$~\cite{NS81}
(see also~\cite[Chapter~7]{LP16}); it is positive because $p>p_c$, and it is not one because
$\PP_p(N_\infty=1)<1$ for $p<p_u$ by the definition of~$p_u$. H\"aggstr\"om, Peres and Schonmann proved
the stronger statement that this holds simultaneously for all $p\in(p_c,p_u)$ in the standard
coupling~\cite{HPS99}. The compact interval in \cref{thm:main} is obtained directly from the
boundedness of a single connection operator (\cref{prop:interval}).
\end{remark}

\begin{remark}[Scope]\label{rem:scope}
The manuscript~\cite{Ope26} treats quasi-transitive graphs, in which $\Aut(G)$ has finitely many
orbits, and also allows loops and multiple edges. A transitive graph with multiple edges does not
reduce to a transitive simple graph: subdividing each edge, as in~\cite{Ope26}, produces a simple
graph with two vertex orbits. Such graphs are therefore not covered by these notes.
\end{remark}
